\section{Action-aligned certificates and selection}\label{supp:action99}
\subsection{Proof of Theorem~\ref{thm:action99}}
Taking the infimum of the second inequality in \eqref{eq:support99} and subtracting it from the first gives the gap bound. Substitution of $b=a$ in that inequality shows $L\le Q(a)\le U$, hence $w\ge0$. The minimum in \eqref{eq:gamma99} is at most zero because $a$ is feasible, so $\Gamma\ge0$.

For $z=b-a$ and the supplied $n\in N_A(a)$,
\[
 s^{\mathsf T}z
 =(s-\widehat s+\widehat s+n)^{\mathsf T}z-n^{\mathsf T}z
 \ge -(d+r)\|z\|_2.
\]
The inequality uses $n^{\mathsf T}z\le0$ and the triangle and Cauchy--Schwarz inequalities. Completing the square yields
\[
 -(d+r)\|z\|_2+\frac\mu2\|z\|_2^2
 \ge -\frac{(d+r)^2}{2\mu}.
\]
Taking the minimum over the original feasible set proves \eqref{eq:direction99}. In one dimension, completion of the square centers the minimization at $a-s/\mu$, whose projection onto the interval gives the stated minimizer. This argument includes $a=0$ and $a=A_0$ and does not differentiate a discontinuous actor. A lower numerical bound on the minimum only enlarges $\Gamma$, preserving validity.\qed

\subsection{Proof of the risk-witness lift}
Strong convexity of the stage cost gives
\[
 c(b)\ge c(a)+\nabla c(a)^{\mathsf T}(b-a)+\frac\mu2\|b-a\|_2^2.
\]
For the common feasible risk witness, the dual representation and the continuation support imply
\[
 \rho(V\circ F(b))\ge\int V(F(b,\omega))\,q(d\omega)-\psi(q)
 \ge\int\ell\,dq-\psi(q)+\left(\int G^{\mathsf T}v\,dq\right)^{\!\mathsf T}(b-a).
\]
Adding the inequalities gives \eqref{eq:lift99}. The same $q$ must be feasible for every $b$ in the comparison. If an upper penalty endpoint replaces $\psi(q)$, the intercept decreases and remains valid. Errors in a vector integral must be retained as a set of possible slopes or transferred uniformly to the intercept using the bounded action displacement. Choosing a favorable slope from an interval without this transfer is invalid. The complete policy upper $U$ must separately dominate the full risk envelope, not just the selected lower witness.\qed

\subsection{Dated action losses and own-future continuations}
Suppose first that a feasible complete policy has, at every date and state,
\[
 Q_t(V^*_{t+1};x,\pi_t(x))-V_t^*(x)\le\gamma_t
\]
with nonnegative constants $\gamma_t$. If $V^\pi_{t+1}\le V^*_{t+1}+e_{t+1}$, monotonicity and cash invariance give
\[
 V_t^\pi\le Q_t(V^*_{t+1};x,\pi_t(x))+\beta_te_{t+1}
 \le V_t^*+\gamma_t+\beta_te_{t+1}.
\]
Backward induction with $e_T=0$ proves $V_0^\pi-V_0^*\le\sum_{t<T}B_t\gamma_t$. A certified terminal-controller allowance is inserted at its actual date. No product of tail multipliers is used to propagate a constant.

This statement does not assume that an NBO receives $V^*_{t+1}$ as a free training target. A support built only for its own attained $f_{t+1}$ bounds $Q_t(f_{t+1})$, not $Q_t(V^*_{t+1})$. It must therefore enter the one-sided or centered residual sandwich with a certified continuation discrepancy. The investment construction instead supplies genuine original-$V_1$ supports and a separately bounded terminal controller. These are two legitimate implementations of the same original-optimum obligation; a small own-future greedy gap alone is not that obligation.

\subsection{Proof of Theorem~\ref{thm:selection99}}
Conditional on acquisition, the menu and all its functions are fixed. For any one $h$ and any loss $Z_h\in[0,1]$, Hoeffding's inequality gives
\[
 \Pr\left(\left|n^{-1}\sum_iZ_h(X_i)-\mathbb EZ_h(X)\right|>e\mid\mathcal A\right)
 \le2e^{-2ne^2}.
\]
A union bound over the $M$ candidates, at $e=e_n$, gives a common event of conditional probability at least $1-\delta$. On that event, for every $h$,
\[
 \mathbb E Z_{\widehat h}
 \le\widehat Z_{\widehat h}+e_n
 \le\widehat Z_h+e_n
 \le\mathbb E Z_h+2e_n.
\]
Use $Z_h=R_h$ to obtain \eqref{eq:rejection99}. Use $Z_h=\ell_h$ and the pointwise inequality $R_h\le\ell_h$ to obtain the surrogate assertion. These are separate choices of loss; claiming both events simultaneously requires a corresponding union allocation. The work implication follows by expectation of its verified pointwise cap. For direct work minimization use $Z_h=W_h/W_{\max}$ and multiply the bound by $W_{\max}$. Bounds exceeding the trivial support limit may be intersected with it. Validation and acquisition costs do not disappear from total workload cost because they are absent from the per-context random variable.\qed

A minimal counterexample explains why a global residual objective cannot replace this loss. Let the action set be $[-1,1]$, the true objective $Q(a)=a^2$, and consider two continuation perturbations added to this objective. The first is $e_0(a)=a^2$, with oscillation one and unchanged minimizer zero. The second is $e_1(a)=\eta a$, with oscillation $2\eta<1$ for $0<\eta<1/2$, but its minimizer is $-\eta/2$ and its true regret is $\eta^2/4>0$. Thus even strict oscillation reduction can worsen the delivered action. This example is an explanatory construction, not a diagnosis inferred for each retained investment cell.

\section{Nonuniform and randomized AVaR information}\label{supp:information99}
\subsection{Proof of the nonuniform lower bound}
Let $q,L,t$ be as in Theorem~\ref{thm:nonuniform99} and answer every query with
\[
 f_0(z)=\frac{K}{2L}(z-q)_+^2.
\]
The function is convex, nondecreasing and $K$-Lipschitz. Collect the queried points strictly inside $(q,R)$, insert the two endpoints, and write the resulting nonempty interval lengths as $h_1,\ldots,h_s$, with $s\le N+1$. Let $f_1=f_0$ to the left of $q$ and equal its chord on each of those intervals. Successive chord slopes are nondecreasing and belong to $[0,K]$, including the junction at $q$. Therefore $f_1\in\mathcal C_{K,R}$ and is nondecreasing. The functions agree at every query, and induction on the adaptive transcript proves that a deterministic procedure has the same output on both.

Because the functions are continuous and nondecreasing and $P$ has positive density, their upper-tail risks are $t^{-1}\int_q^R f(z)p(z)\,dz$. Flat portions cause no difficulty: the quantile formula integrates the appropriate fraction of their mass. On an interval of length $h$, the unweighted integral of the chord error is $Kh^3/(12L)$. Since $p\ge m$,
\[
 \operatorname{AVaR}_\alpha^P(f_1)-\operatorname{AVaR}_\alpha^P(f_0)
 \ge\frac{mK}{12Lt}\sum_{i=1}^s h_i^3
 \ge\frac{mKL^2}{12t(N+1)^2}.
\]
The last inequality follows from $\sum_i h_i=L$ and convexity of the cube. A sound common interval contains both values. Finally $t=\int_q^R p\le M_pL$, so $L\ge t/M_p$, giving the second lower bound. The worst-case error of a point-valued deterministic procedure is at least half the first width.\qed

\subsection{Proof of the nonuniform upper bound}
Write the uniform cell length as $h=2R/n$. For a cell $[a,b]$, let its mass be $p_i$ and its conditional mean be $\mu_i$. For any convex function $g$,
\[
 p_i g(\mu_i)\le\int_{a}^{b}g(z)p(z)\,dz
 \le p_i\left\{\frac{b-\mu_i}{h}g(a)+\frac{\mu_i-a}{h}g(b)\right\}.
\]
For $g_\eta=(f-\eta)_+$, these are valid for every threshold $\eta$; $g_\eta$ is convex and $K$-Lipschitz. If $a<\mu_i<b$, the difference of the right and left expressions is
\[
 p_i\frac{(\mu_i-a)(b-\mu_i)}{h}
 \left\{\frac{g_\eta(b)-g_\eta(\mu_i)}{b-\mu_i}
          -\frac{g_\eta(\mu_i)-g_\eta(a)}{\mu_i-a}\right\}.
\]
The factor before the braces is at most $p_i h/4\le M_ph^2/4$. The slope difference is bounded by the variation of the one-sided derivative within the cell. Summing these nonnegative variations over the disjoint cell interiors is at most the total derivative variation, which is at most $2K$. Boundary means follow by continuity (and do not occur with strictly positive density). Hence the summed upper-minus-lower expectation is at most $M_pKh^2/2=2M_pKR^2/n^2$, uniformly in $\eta$.

Add $\eta$, divide each expectation by $t$, and minimize over $\eta$ in the AVaR variational representation. The numerical lower objective lies below the original one and the numerical upper objective above it. Their difference is uniformly bounded by \eqref{eq:upper99}; taking infima preserves both enclosure and this width bound. The numerical rules require the $n$ mean values and the $n+1$ endpoint values of $f$, with repetitions shared. Known moments specify their nonnegative weights, summing to one. An approximate moment rule requires a separately proved perturbation allowance and is not covered by an uncharged exact-law oracle.\qed

\subsection{Proof for pathwise-sound random intervals}
Use the same fixed $f_0$ as above, before drawing a random tape. Outside the single null set in the soundness definition, fixing the tape gives a deterministic adaptive procedure that is sound for all inputs. The preceding chord construction therefore forces the stated width on its transcript from $f_0$. The chord competitor may vary with the tape; $f_0$ does not. Taking expectation gives the same lower bound on expected width. A procedure whose bad-tape set depends on the input need not be sound for that tape's chord competitor, which is why the argument proves no such assertion for ordinary $1-\delta$ coverage.\qed

\subsection{Proof of the randomized point-estimation lower bound}
For the uniform law set $q=-R+2R\alpha$, $L=2R(1-\alpha)$ and $s=z-q$. Partition the tail into $2N$ equal intervals of length $h=L/(2N)$. Let $\psi(u)=u^2(1-u)^2$ on $[0,1]$, and extend it by zero outside this interval. For independent signs $\theta_j\in\{-1,1\}$, define on the tail
\[
 f_\theta(q+s)=\frac{K s^2}{4L}
 +\frac{Kh^2}{8L}\sum_{j=1}^{2N}\theta_j
       \psi\left(\frac{s-(j-1)h}{h}\right),
\]
and put $f_\theta=0$ to the left of $q$. We verify that this entire finite family belongs to the announced class. Both $\psi$ and $\psi'$ vanish at zero and one, so the functions and first derivatives match at every cell boundary and at $q$. On each cell, $\psi''(u)=2-12u+12u^2\in[-1,2]$. Consequently
\[
 f_\theta''\ge\frac K{2L}-\frac K{8L}\,2=\frac K{4L}>0.
\]
The derivative starts at zero and ends at $K/2$, since bump derivatives vanish at both tail endpoints. It is therefore everywhere in $[0,K/2]$. The full function is convex, nondecreasing and $K$-Lipschitz; second derivatives at junctions are not needed.

Since $\int_0^1\psi(u)\,du=1/30$, the original AVaR of this monotone function equals a known baseline plus
\[
 b_N\sum_{j=1}^{2N}\theta_j,
 \qquad b_N=\frac{Kh^3}{240L^2}=\frac{KL}{1920N^3}.
\]
A point query inside one open cell depends only on that cell's sign; a boundary query or query outside the tail reveals none. Conditional on the algorithm's random tape and its entire adaptive transcript, the signs in unqueried cell interiors remain independent and equally likely. This follows inductively: the next location is measurable with respect to already revealed information and the tape, not unqueried signs. At most $N$ of the $2N$ signs have been queried. The conditional variance of the target risk is therefore at least $Nb_N^2$.

For any estimator measurable with respect to the transcript and independent algorithmic randomization, conditional mean-square error is no smaller than that conditional variance. Averaging over signs and randomization gives Bayes risk at least $Nb_N^2$. The largest risk over this fixed finite family is at least its average; hence the worst-case root mean-square error is at least
\[
 \sqrt N b_N=\frac{KL}{1920N^{5/2}}
             =\frac{KR(1-\alpha)}{960N^{5/2}}.
\]
This argument uses a fixed distribution over inputs rather than a tape-dependent adversarial input, and therefore applies to ordinary randomized point estimation. It supplies a lower bound only; an optimal upper algorithm for the whole convex class is not asserted.\qed

\subsection{Proof of the dated allocation identity}
For positive $a_j,c_j$, put $H=\sum_j a_j^{1/3}c_j^{2/3}$. H\"older's inequality gives
\[
 H=\sum_j(a_j/n_j^2)^{1/3}(c_jn_j)^{2/3}
 \le\left(\sum_j a_j/n_j^2\right)^{1/3}
       \left(\sum_j c_jn_j\right)^{2/3}.
\]
Thus feasibility implies $\sum_jc_jn_j\ge H^{3/2}/\sqrt\varepsilon$. Equality is attained by
$n_j=(a_j/c_j)^{1/3}\sqrt{H/\varepsilon}$, for which the tolerance constraint binds. Upward integer rounding decreases each allowance and adds less than $c_j$ per date. This argument assumes the stated linear per-resolution work law on a fixed graph; it does not substitute for pricing descendants or multidimensional tensor queries.\qed

\section{Decision separation and exact acquisition}\label{supp:decision99}
\subsection{Proof of Theorem~\ref{thm:decision99}}
At fixed $x$, only prices between $h(x)$ and $v(x)$ induce a different action from the original optimum. At such a price, the regret is $|p-v(x)|$. Both thresholds lie in $[0,3]$, so integrating the triangle under the uniform price density gives
\[
 \frac13\int_{\min(h,v)}^{\max(h,v)}|p-v|\,dp
 =\frac{(h-v)^2}{6}.
\]
Integration over $x$ proves the unconditional identity, with no approximation or simulated policy observation.

Write $f=v-1=(x_1+x_2-1)_+$. Under independent uniform coordinates,
\[
 \mathbb Ef=\frac16,\quad
 \mathbb E[f\mid x_1]=\frac{x_1^2}{2},\quad
 \mathbb E[f\mid x_2]=\frac{x_2^2}{2},\quad
 \mathbb Ef^2=\frac1{12}.
\]
The additive projection of $f$ is $g=(x_1^2+x_2^2)/2-1/6$. Its residual has conditional mean zero given either coordinate, hence is orthogonal to every square-integrable additive function. This proves optimality over the whole additive class, not merely over additive quadratics. The constant and centered coordinate terms are mutually orthogonal, so
\[
 \|g\|_2^2=\frac1{36}+2\left(\frac1{20}-\frac1{36}\right)=\frac{13}{180},
 \qquad \|f-g\|_2^2=\frac1{12}-\frac{13}{180}=\frac1{90}.
\]
The shifted optimum $1+g$ has range $[5/6,11/6]$, satisfying the price-support restriction. Equation~\eqref{eq:decisionloss99} gives $1/540$.

For the quadratic projection use the orthogonal shifted Legendre polynomials
$P_0=1$, $P_1(u)=2u-1$, $P_2(u)=6u^2-6u+1$, with squared norms $1/(2j+1)$. Direct integration on the upper triangle gives
\[
 \Pi_2v=\frac76+\frac14(P_1(x_1)+P_1(x_2))
 +\frac1{12}(P_2(x_1)+P_2(x_2))
 +\frac3{10}P_1(x_1)P_1(x_2).
\]
To verify these coefficients without assuming the answer, the monomial normal equations use basis $(1,x_1,x_2,x_1^2,x_1x_2,x_2^2)$, Gram entries
$G_{(i,j),(k,l)}=1/\{(i+k+1)(j+l+1)\}$, and right side
\[
 b_{ij}=\frac1{(i+1)(j+1)}+
 \int_0^1\int_{1-x_1}^1x_1^ix_2^j(x_1+x_2-1)\,dx_2\,dx_1.
\]
The Gram matrix is positive definite because a nonzero polynomial cannot vanish almost everywhere on the square. Substituting the coefficients in \eqref{eq:quadratic99} gives $G\theta=b$ exactly; the complete rational matrix, vector and zero residual are supplied in the executable evidence. Orthogonality alternatively gives
\[
 \|\Pi_2 f\|_2^2
 =\frac1{36}+2\frac{(1/4)^2}{3}
       +2\frac{(1/12)^2}{5}+\frac{(3/10)^2}{9}
 =\frac{37}{450},
 \quad \|f-\Pi_2 f\|_2^2=\frac1{900}.
\]
The Hessian of \eqref{eq:quadratic99} has eigenvalues $11/5$ and $-1/5$. Its interior critical point $(3/11,3/11)$ is a saddle. Checking each edge gives the minimum $143/150$ at $(3/5,0)$ or $(0,3/5)$ and the maximum $32/15$ at $(1,1)$. Thus it also lies in $[0,3]$, and the constrained admissibility requirement does not change the projection optimum. Dividing by six gives $1/5400$.

The original $v$ is the displayed affine constant plus a single ReLU, and has zero error. Its native hinge, two-triangle piecewise-affine encoding, and dictionary containing that hinge are exactly the same function, so they induce the same policy at every state and price. No representational renaming changes either conclusion.\qed

For the worst-state conditional criterion, \eqref{eq:decisionloss99} gives
$\sup_x\mathbb E_p[J^{b_h}-J^*]=\|h-v\|_\infty^2/6$. The sharp attained approximation errors in Proposition~\ref{prop:ridge97} therefore yield $1/96$ and $1/1536$. The explicit attaining additive and quadratic functions shifted by one stay inside $[0,3]$. The population and worst-state optimization problems are distinct, so their minimizers need not coincide. The exact purchase equation follows by subtracting the respective policy and complete computation costs and solving the strict linear inequality.\qed
